\documentclass[11pt,a4paper]{article}
\usepackage[margin=1in]{geometry}
\usepackage[utf8]{inputenc}
\usepackage[T1]{fontenc}
\usepackage{lmodern}
\usepackage{hyperref}
\usepackage{enumitem}
\usepackage{booktabs}
\usepackage{xcolor}

\hypersetup{
    colorlinks=true,
    linkcolor=blue,
    citecolor=blue,
    urlcolor=blue
}

\title{\textbf{Data Availability Statement} \\[0.5em]
{\large Target breadth and mutation resilience in African-natural-product-inspired antimalarial chemotypes: a computational analysis}}

\author{
Myke Vital Sao Temgoua, Jean-Pierre Tchapet Njafa, Penabei Samafou, \\
Wilfred Fon Mbacham, Serge Guy Nana Engo
}

\date{\today}

\begin{document}

\maketitle

\section*{Summary}

All data, code, computational protocols, and materials necessary to reproduce the results presented in this manuscript are publicly available through a permanent Zenodo repository under an open-source MIT license. This ensures complete computational reproducibility and supports open science practices.

\section*{Primary Data Repository}

\subsection*{Zenodo Deposit}

\begin{itemize}[leftmargin=*]
\item \textbf{Repository:} Zenodo (CERN-hosted permanent repository)
\item \textbf{DOI:} \url{https://doi.org/10.5281/zenodo.22696778}
\item \textbf{License:} MIT License (Open Source Initiative approved)
\item \textbf{Accessibility:} Publicly accessible without registration or fees
\item \textbf{Preservation:} Permanent archival with version control
\item \textbf{Verification:} SHA256 checksums for all files
\end{itemize}

\section*{Repository Contents}

The Zenodo deposit contains the following computational materials organized in a versioned directory structure:

\subsection*{1. Source Code and Scripts}

\begin{itemize}[leftmargin=*]
\item \textbf{Chemical-space generation:} Python scripts for hybrid library construction, SMILES variational autoencoder, K-means clustering, centroid selection, and novelty analysis
\item \textbf{Docking workflows:} AutoDock Vina configuration files, grid specifications, Meeko ligand preparation scripts, and geometric quality-control filters
\item \textbf{RRS analysis:} Per-target resistance-resilience score computation, classification algorithms, and wild-type binding-threshold filters
\item \textbf{Statistical analysis:} Spearman correlation scripts, Bonferroni correction, permutation tests, power analysis calculations
\item \textbf{Validation protocols:} DEKOIS 2.0 enrichment analysis, MMV Malaria Box screening, redocking RMSD calculations, retrospective antimalarial panel evaluation
\item \textbf{Visualization:} Figure generation scripts for chemical-space coverage, RRS profiles, cross-metric relationships, and target-wise distributions
\end{itemize}

\subsection*{2. Input Data}

\begin{itemize}[leftmargin=*]
\item \textbf{Seed libraries:} 396 African natural product SMILES, 454 synthetic antimalarial SMILES with source annotations
\item \textbf{Receptor structures:} Prepared PDB files for PfDHFR (7F3Y), PfCRT (6UKJ corrected to LYS-76), PfClpP (2F6I), and PfATP4 (9N10)
\item \textbf{Ligand structures:} SMILES strings and 3D coordinates for all 17 Set-C candidates in canonical and protonation states
\item \textbf{Mutation definitions:} Receptor coordinate files for PfDHFR mutants (N51I, C59R, S108N, I164L) and PfCRT mutants (K76T, K76A)
\item \textbf{Validation datasets:} DEKOIS 2.0 PfDHFR panel (40 actives, 1200 decoys), MMV Malaria Box subset, reference ligand coordinates for redocking
\item \textbf{Retrospective panel:} Structures for artemisinin, pyrimethamine, chloroquine, lumefantrine, and cipargamin
\end{itemize}

\subsection*{3. Computational Protocols}

\begin{itemize}[leftmargin=*]
\item \textbf{Grid specifications:} Complete AutoDock Vina grid box centers, dimensions, and exhaustiveness parameters for all four targets
\item \textbf{Docking parameters:} Meeko ligand preparation settings, protonation-state assignment rules, PDBQT conversion protocols
\item \textbf{Quality gates:} Geometric pose-quality criteria (centroid position, anchor distance, box containment)
\item \textbf{RRS definitions:} Wild-type binding thresholds, eligibility criteria, classification boundaries (A*, A, B, C, D)
\item \textbf{Multi-seed protocol:} Seed values, exhaustiveness settings, and convergence criteria for sensitivity analysis
\item \textbf{Null controls:} Pipeline-null protocol for PfCRT K76T/K76A signal assessment
\end{itemize}

\subsection*{4. Raw Computational Outputs}

\begin{itemize}[leftmargin=*]
\item \textbf{Docking poses:} Complete PDBQT files for all candidate--target pairs (68 wild-type, 102 mutant-panel)
\item \textbf{Vina score logs:} Raw scoring-function outputs for every docking run with mode ranks and RMSD values
\item \textbf{Cluster assignments:} K-means cluster labels and centroid coordinates for the 65,856-molecule library
\item \textbf{Fingerprint matrices:} Morgan fingerprint bit vectors for novelty and scaffold-similarity calculations
\item \textbf{RRS matrices:} Per-candidate, per-target, per-mutant score ratios and classification assignments
\item \textbf{Validation metrics:} ROC curves, AUC values, enrichment factors, redocking RMSD distributions
\end{itemize}

\subsection*{5. Machine-Readable Tables}

\begin{itemize}[leftmargin=*]
\item \textbf{Candidate manifest:} Complete Set-C roster with SMILES, identifiers, MPO scores, and source annotations (CSV, JSON)
\item \textbf{Vina matrix:} $17 \times 4$ scoring table with target-specific wild-type estimates (CSV, JSON)
\item \textbf{RRS profiles:} Per-candidate mutant-panel results with eligibility flags and class assignments (CSV, JSON)
\item \textbf{Cross-metric data:} PNS, ACSI, wild-type scores, favorable-target counts, and correlation statistics (CSV, JSON)
\item \textbf{Physicochemical properties:} Lipinski/Veber descriptors, QED values, and ADMET predictions (CSV, JSON)
\item \textbf{Validation datasets:} DEKOIS actives/decoys, MMV screening results, redocking reference coordinates (CSV, JSON)
\end{itemize}

\subsection*{6. Documentation}

\begin{itemize}[leftmargin=*]
\item \textbf{README files:} Installation instructions, software dependencies, execution workflows
\item \textbf{Data dictionary:} Column definitions, unit specifications, controlled vocabularies
\item \textbf{Provenance records:} Software versions (Python, RDKit, AutoDock Vina, Open Babel), execution timestamps, hardware specifications
\item \textbf{File manifest:} SHA256 checksums for all deposited files with verification scripts
\item \textbf{License file:} Full text of MIT License covering all repository contents
\end{itemize}

\section*{Software and Dependencies}

All computational work was performed using open-source software. The specific versions used are documented in the Zenodo repository:

\begin{itemize}[leftmargin=*]
\item \textbf{Python:} 3.11 (conda environment specification included)
\item \textbf{RDKit:} 2025.03.6 (cheminformatics toolkit)
\item \textbf{AutoDock Vina:} 1.2.7 (molecular docking)
\item \textbf{Meeko:} 0.5.0 (ligand preparation)
\item \textbf{Open Babel:} 3.1.1 (molecular format conversion)
\item \textbf{NumPy, Pandas, Matplotlib, Seaborn:} Standard scientific Python stack (versions in requirements.txt)
\item \textbf{scikit-learn:} Clustering and machine learning utilities
\item \textbf{SELFIES:} 2.1.1 (molecular string representation)
\end{itemize}

\section*{Reproducibility Instructions}

The Zenodo repository includes complete step-by-step instructions for:

\begin{enumerate}[leftmargin=*]
\item Environment setup (conda YAML file provided)
\item Downloading and verifying all input files (SHA256 checksums)
\item Executing the complete workflow from seed libraries to final results
\item Regenerating all manuscript figures and tables
\item Running validation protocols (DEKOIS, MMV, redocking)
\item Reproducing statistical analyses and cross-metric correlations
\end{enumerate}

\section*{External Database Sources}

The following public databases were used to construct the seed libraries and validation sets. All accessions and retrieval dates are documented:

\begin{itemize}[leftmargin=*]
\item \textbf{ZINC15:} Commercially available compound library (\url{https://zinc15.docking.org/})
\item \textbf{ENAMINE-REAL:} Synthesizable compound space (\url{https://enamine.net/compound-libraries/real-compounds})
\item \textbf{African Natural Products Database (AfroDB):} Curated natural product collection
\item \textbf{RCSB Protein Data Bank:} Target structures 7F3Y, 6UKJ, 2F6I, 9N10 (\url{https://www.rcsb.org/})
\item \textbf{DEKOIS 2.0:} Demanding evaluation kits for objective in silico screening (\url{http://www.dekois.com/})
\item \textbf{MMV Malaria Box:} Medicines for Malaria Venture compound collection (\url{https://www.mmv.org/})
\item \textbf{ChEMBL:} Bioactivity database for approved antimalarials (\url{https://www.ebi.ac.uk/chembl/})
\end{itemize}

\section*{Data Usage and Citation}

All data and code in the Zenodo repository are released under the MIT License, which permits:

\begin{itemize}[leftmargin=*]
\item Free use, modification, and distribution
\item Commercial and non-commercial applications
\item Derivative works with attribution
\end{itemize}

Users of these materials should cite:

\begin{quote}
Sao Temgoua, M.V.; Tchapet Njafa, J.-P.; Samafou, P.; Mbacham, W.F.; Nana Engo, S.G. \textit{Target breadth and mutation resilience in African-natural-product-inspired antimalarial chemotypes: a computational analysis}. \emph{ChemRxiv}, 2026. DOI: \url{https://doi.org/10.26434/chemrxiv.15006437/v4}

Data and code repository: \url{https://doi.org/10.5281/zenodo.22696778}
\end{quote}

\section*{Contact for Data Inquiries}

For questions regarding data access, interpretation, or technical issues with the repository:

\begin{itemize}[leftmargin=*]
\item \textbf{Corresponding Author:} Myke Vital Sao Temgoua
\item \textbf{Email:} \href{mailto:myke-vital.sao@facsciences-uy1.cm}{myke-vital.sao@facsciences-uy1.cm}
\item \textbf{Institution:} Department of Physics, University of Yaoundé I, Cameroon
\end{itemize}

\section*{Long-Term Preservation}

\begin{itemize}[leftmargin=*]
\item \textbf{Zenodo commitment:} CERN guarantees data preservation for the lifetime of CERN (minimum 20+ years)
\item \textbf{Redundancy:} Multiple backup copies maintained by CERN infrastructure
\item \textbf{Format stability:} All files use non-proprietary, widely-supported formats (CSV, JSON, PDB, PDBQT, TXT, PDF)
\item \textbf{Version control:} Immutable DOI ensures consistent access to the exact version used in this publication
\item \textbf{Migration plan:} CERN commits to format migration if technology changes
\end{itemize}

\section*{Compliance and Ethics}

\begin{itemize}[leftmargin=*]
\item \textbf{No human subjects:} This is a computational study with no human participants
\item \textbf{No animal experiments:} All work is in silico
\item \textbf{No biological materials:} No wet-lab experiments were performed
\item \textbf{Public data only:} All seed compounds and protein structures are from public databases
\item \textbf{No proprietary data:} All computational results are original work by the authors
\item \textbf{Open science principles:} Full transparency and reproducibility prioritized
\end{itemize}

\section*{Acknowledgment}

The authors acknowledge Zenodo and CERN for providing permanent, open-access data hosting infrastructure that enables reproducible computational research. We thank the open-source software communities (RDKit, AutoDock, Python scientific stack) whose tools made this work possible.

\vspace{1em}

\noindent\rule{\textwidth}{0.4pt}

\begin{center}
\textbf{For immediate access to all data and code:} \\[0.5em]
\Large\url{https://doi.org/10.5281/zenodo.22696778}
\end{center}

\noindent\rule{\textwidth}{0.4pt}

\end{document}
